\documentclass[11pt, a4paper]{article}
\usepackage[utf8]{inputenc}
\usepackage{amsmath, amssymb}
\usepackage{graphicx}
\usepackage[margin=1in]{geometry}
\usepackage{tcolorbox}
\usepackage{fancyhdr}
\usepackage{hyperref}
\usepackage{listings}
\usepackage{xcolor}

\definecolor{UNIBBlue}{RGB}{0, 51, 102}
\definecolor{codegreen}{rgb}{0,0.6,0}
\definecolor{codegray}{rgb}{0.5,0.5,0.5}
\definecolor{codepurple}{rgb}{0.58,0,0.82}
\definecolor{backcolour}{rgb}{0.95,0.95,0.92}

\lstdefinestyle{mystyle}{
    backgroundcolor=\color{backcolour},   
    commentstyle=\color{codegreen},
    keywordstyle=\color{magenta},
    numberstyle=\tiny\color{codegray},
    stringstyle=\color{codepurple},
    basicstyle=\ttfamily\footnotesize,
    breakatwhitespace=false,         
    breaklines=true,                 
    captionpos=b,                    
    keepspaces=true,                 
    numbers=left,                    
    numbersep=5pt,                  
    showspaces=false,                
    showstringspaces=false,
    showtabs=false,                  
    tabsize=2
}
\lstset{style=mystyle}

\pagestyle{fancy}
\fancyhf{}
\rhead{\textbf{Optimisasi untuk Teknik Elektro}}
\lhead{Lembar Kerja Mahasiswa (LKM) - Minggu 4}
\cfoot{\thepage}

\begin{document}

\begin{center}
    \Large\color{UNIBBlue}\textbf{LEMBAR KERJA MAHASISWA (LKM)}\\
    \large\textbf{Minggu 4: Membuat dan Menginterpretasikan Laporan Sensitivitas}
\end{center}

\vspace{0.5cm}
\textbf{Nama Praktikan :} \rule{6cm}{0.4pt} \\
\textbf{NPM :} \rule{6.8cm}{0.4pt}

\section*{Tujuan Praktikum}
1. Mahasiswa mampu menghasilkan laporan rentang sensitivitas menggunakan Julia/JuMP atau \textit{solver} di Python. \\
2. Mahasiswa mampu menerjemahkan keluaran komputasional ke dalam strategi operasi keteknikan.

\section*{Kasus: Alokasi Produksi Baterai (Dimensi C3 - C4)}
Sebuah pabrik memproduksi baterai Li-Ion untuk Drone ($x_1$) dan Baterai untuk Mobil Listrik EV ($x_2$). Profit untuk setiap batch baterai Drone adalah Rp 300 Juta, dan EV adalah Rp 500 Juta. Pabrik dibatasi oleh kapasitas mesin Assembly (A) dan mesin Testing (T).
\begin{itemize}
    \item Perakitan: 2 jam/$x_1$ dan 4 jam/$x_2$. Maksimal tersedia 160 jam. ($2x_1 + 4x_2 \le 160$)
    \item Pengujian: 3 jam/$x_1$ dan 2 jam/$x_2$. Maksimal tersedia 120 jam. ($3x_1 + 2x_2 \le 120$)
\end{itemize}

\section*{Aktivitas 1: Ekstraksi Rentang Optimasi (C3)}
Pada \textit{solver} LP modern (seperti Linprog pada SciPy atau opsi tertentu di HiGHS), kita dapat mengekstrak sensitivitas rentang objektif dan RHS.
\textbf{Implementasi Python (SciPy):}
\begin{lstlisting}[language=Python]
from scipy.optimize import linprog

# C = [-300, -500] untuk maksimasi
c = [-300, -500]
A = [[2, 4], [3, 2]]
b = [160, 120]

# Metode highs mendukung analisis margin
res = linprog(c, A_ub=A, b_ub=b, method='highs')

print("Solusi:", res.x)
print("Profit Maksimal:", -res.fun)
print("Shadow Prices (Inequalities):", -res.ineqlin.marginals)
\end{lstlisting}
\textit{Catatan: Di kelas, kita akan menggunakan \textit{package} tambahan atau perhitungan manual sederhana dari matriks basis optimal jika menggunakan bahasa Julia standar, namun konsepnya tetap sama.}

\section*{Aktivitas 2: Analisis Skenario Berdasarkan Rentang (C4)}
Anggap laporan analisis yang dihasilkan sistem menyatakan untuk Baterai Drone ($c_1 = 300$):
\begin{itemize}
    \item \textbf{Allowable Increase}: 50
    \item \textbf{Allowable Decrease}: 112.5
\end{itemize}
Untuk Mesin Assembly ($b_1 = 160$):
\begin{itemize}
    \item \textbf{Shadow Price}: 112.5
    \item \textbf{Allowable Increase}: 80
    \item \textbf{Allowable Decrease}: 80
\end{itemize}

\textbf{Pertanyaan:}
1. Jika terjadi persaingan pasar yang memaksa kita menurunkan profit Baterai Drone menjadi Rp 250 Juta/batch, apakah strategi alokasi produksi kita saat ini harus diubah? Jelaskan!
2. Jika ada tawaran sewa mesin perakitan tambahan seharga Rp 100 Juta/jam dengan kuota 50 jam, apakah tawaran ini menguntungkan untuk diambil? Evaluasi menggunakan prinsip Range of Feasibility dan Shadow Price.

\vspace{1cm}
\textbf{Jawaban:} \\
\rule{\textwidth}{0.4pt} \\
\rule{\textwidth}{0.4pt} \\
\rule{\textwidth}{0.4pt} \\
\rule{\textwidth}{0.4pt}

\end{document}
